\section*{Chapitre \ref{ch:ll:where-latency-comes-from} --- D'où vient la latence}

\begin{solution}{exo:ll:where-latency-comes-from:1}
$45 \times 1.462 / (2.998\times10^8) = \qty{219}{\nano\second}$ dans chaque sens.
\end{solution}

\begin{solution}{exo:ll:where-latency-comes-from:2}
$148 \times 8 = 1\,184$ bits~: \qty{118.4}{\nano\second} à 10~Gb/s et \qty{47.4}{\nano\second} à 25~Gb/s.
\end{solution}

\begin{solution}{exo:ll:where-latency-comes-from:3}
Les moyennes s'additionnent~: \qty{1\,030}{\nano\second}. Rien d'exact ne se dit de la médiane de la somme à partir des moyennes ou
médianes~; pour des étapes indépendantes à queue légère, elle est en général proche de la somme des médianes, comme dans le chemin synthétique du chapitre
(\qty{2.17}{\micro\second} contre \qty{2.13}{\micro\second}).
\end{solution}

\begin{solution}{exo:ll:where-latency-comes-from:4}
$1 - 0.996^5 = 1.98~\%$. Comme plus de 1~\% des messages rencontrent un blocage, le p99 du chemin se situe dans les blocages, bien qu'aucun
p99 d'étape ne s'y trouve.
\end{solution}

\begin{solution}{exo:ll:where-latency-comes-from:5}
En appliquant la \cref{prop:ll:where-latency-comes-from:union} avec $\alpha_i = 0.001/4$~: chaque étape au 99.975\up{e}~percentile, avec
des cibles dont la somme ne dépasse pas \qty{10}{\micro\second}.
\end{solution}

\begin{solution}{exo:ll:where-latency-comes-from:6}
À \qty{2.0}{\micro\second}, supprimer les blocages fait passer $\P(\text{win})$ de 0.662 à 0.675 (1.3 points). Sur la
courbe avec blocages, $\P(\text{win})$ augmente d'environ 0.81 par microseconde au voisinage de \qty{2}{\micro\second}, si bien que le même gain
fait baisser la médiane d'environ \qty{16}{\nano\second}.
\end{solution}

\begin{solution}{exo:ll:where-latency-comes-from:7}
\texttt{ll\_budget.composition(stall\_p=0.05)}~: le p99 de chaque étape se situe désormais dans ses blocages, et la somme des six
p99 d'étape vaut \qty{79}{\micro\second}, contre un p99 de bout en bout de \qty{32}{\micro\second}. La comparaison s'est inversée~: la
somme des percentiles sous-estime lorsque les blocages sont plus rares que le percentile, et la surestime lorsqu'ils sont fréquents et ne
coïncident pas. Ni l'une ni l'autre n'est un budget.
\end{solution}

\begin{solution}{exo:ll:where-latency-comes-from:8}
Les percentiles ne s'additionnent pas. Les trois p99 ont été mesurés lors d'exécutions séparées, si bien qu'ils ne disent rien de la manière dont les messages lents des étapes
se recoupent~: des blocages rares que chaque banc d'essai exclut de son p99 peuvent placer le p99 du pipeline bien au-dessus de
\qty{1.3}{\micro\second} (\cref{ex:ll:where-latency-comes-from:nonsub}), et des blocages communs aux trois peuvent le placer
en dessous. Mesurer le pipeline de bout en bout, message par message, avec des horodatages à chaque frontière d'étape.
\end{solution}

\begin{solution}{pb:ll:where-latency-comes-from:1}
\begin{enumerate} \sloppy
\item $\P(X_A < X_B)$, les ex æquo étant partagés~; elle suppose que notre \omterm{def:ll:where-latency-comes-from:latency}{latence} et celle du concurrent sur un message donné sont
indépendantes (une rafale qui ralentit les deux invalide cette hypothèse).
\item $\ln X_A - \ln X_B \sim \mathcal N(\ln(2.0/2.2), 2\times0.15^2)$, donc $\P(\text{win}) =
\Phi\big(\ln(2.2/2.0)/(0.15\sqrt2)\big) = \Phi(0.449) = 0.673$.
\item 66.1~\%.
\item Nous ne perdons que les courses bloquées que nous aurions autrement gagnées, soit environ $0.02 \times 0.67$~; et les blocages propres au concurrent
nous en font gagner quelques-unes.
\item p50 \qty{2.01}{\micro\second}~; p99 \qty{23.6}{\micro\second}.
\item Diviser la médiane par deux~: 98.0~\%. Diviser le taux de blocage par deux~: 66.7~\%, le p99 tombant à \qty{3.5}{\micro\second}.
\item 67.4~\%, soit 1.3 points de plus.
\item La médiane.
\item Les blocages débutent à $10^5 \times 0.02 = 2\,000$ par seconde et durent \qty{30}{\micro\second}~: 6~\% du temps.
\item $20 \times 0.06 = 1.2$ exécutions périmées par seconde à $0.5 \times \$0.01 \times 100 = \$0.50$~: 0.60~\$ par seconde,
14\,040~\$ par jour.
\item Elle divise ce montant par deux, soit 7\,020~\$ par jour.
\item L'exposition dépend de la queue (taux de blocage multiplié par la durée des blocages)~; la course, elle, dépend du corps de la distribution.
\item Environ \qty{1.64}{\micro\second}.
\item $1.64 - 1.0 = \qty{0.64}{\micro\second}$ réparties sur cinq étapes~: environ \qty{128}{\nano\second} chacune.
\item Au 99.8\up{e}~percentile ($1 - 0.01/5$).
\item \textbf{Résultat à retenir.} Nous gagnons 66.1~\% des courses~; diviser notre médiane par deux porte ce taux à 98.0~\%, alors que diviser notre taux de blocage par deux ne le fait monter
qu'à 66.7~\%, bien qu'elle ramène notre p99 de \qty{23.6}{\micro\second} à \qty{3.5}{\micro\second}.
\item Les courses dépendent de tout le corps des deux distributions, et l'exposition de la queue~; un seul nombre ne résume ni l'une ni l'autre.
\item Sa \omterm{def:ll:where-latency-comes-from:latency}{latence} n'est visible qu'à travers ses résultats~: les horodatages de ses ordres comparés aux nôtres dans les
messages de la bourse, ou la part des courses perdues en fonction de notre propre \omterm{def:ll:where-latency-comes-from:latency}{latence} mesurée.
\item Un algorithme d'exécution piloté par un calendrier, qui écoule un ordre de grande taille sur plusieurs heures.
\item Une définition partagée et réfutable de «~assez rapide~» pour chaque étape, afin que chaque changement soit mesuré à l'aune de celle-ci.
\end{enumerate}
\end{solution}

\begin{solution}{iq:ll:where-latency-comes-from:1}
La \omterm{def:ll:where-latency-comes-from:t2t}{latence tick-to-trade} couvre l'intervalle allant du message de données de marché déclencheur à l'ordre qui en résulte~; la \omterm{def:ll:where-latency-comes-from:t2t}{latence wire-to-wire} est le même
intervalle mesuré sur le câble, par un appareil de capture doté d'horodatages matériels placé sur une prise de dérivation ou un commutateur, si bien qu'elle inclut les cartes réseau
et le système d'exploitation. Dans le programme, horodater l'arrivée du paquet et la transmission de l'ordre avec
l'horloge du processeur.
\iqlookfor{les deux instants nommés explicitement, et qui prend chaque horodatage.}
\end{solution}

\begin{solution}{iq:ll:where-latency-comes-from:2}
Les pertes se concentrent dans les messages lents (cotations périmées, courses manquées), et la moyenne les masque~; les percentiles élevés montrent
la queue, et la médiane le message ordinaire.
\iqlookfor{que la \omterm{def:ll:where-latency-comes-from:latency}{latence} est une distribution et que l'argent se perd dans sa queue.}
\end{solution}

\begin{solution}{iq:ll:where-latency-comes-from:3}
Pas grand-chose~: les percentiles ne s'additionnent pas. La borne de Boole donne $q_{0.97} \le \qty{6}{\micro\second}$ quelle que soit la dépendance, mais
le p99 du pipeline peut dépasser \qty{6}{\micro\second} si chaque composant présente des blocages rares sous son propre niveau de 1~\%. Mesurer le
pipeline de bout en bout.
\iqlookfor{la non-additivité des quantiles et la borne de Boole.}
\end{solution}

\begin{solution}{iq:ll:where-latency-comes-from:4}
Les chemins de réception et d'émission de la carte réseau, les interruptions et la pile réseau du noyau (ou une bibliothèque de
contournement du noyau), la copie, le décodage, la mise à jour du carnet, la décision, les contrôles de risque, l'encodage et l'attente~: que le fil soit
ordonnancé, pour un verrou, pour une mémoire qui n'est pas dans le cache.
\iqlookfor{le matériel, le système d'exploitation et les étapes logicielles, et les blocages qui les séparent.}
\end{solution}

\begin{solution}{iq:ll:where-latency-comes-from:5}
Les courses se jouent sur le corps des deux distributions~: à 60~\%, la médiane est là où se trouve le rendement, sauf si la
mesure montre que les courses perdues coïncident avec nos blocages. La queue mérite un investissement pour une autre raison~: l'exposition des cotations au repos
tant qu'elles sont périmées.
\iqlookfor{la distinction entre les issues des courses et l'exposition, et une demande de données avant toute dépense.}
\end{solution}

\begin{solution}{iq:ll:where-latency-comes-from:6}
Capturer les paquets entrants et sortants sur une prise de dérivation dotée d'horodatages matériels pour la mesure wire-to-wire~; ajouter des horodatages avec
l'horloge du processeur à chaque frontière d'étape dans le programme, enregistrés hors du \omterm{def:ll:where-latency-comes-from:hot}{chemin critique} (hot path)~; joindre les deux par
identifiants de message~; rejouer une journée chargée enregistrée à son rythme d'origine~; présenter par étape et de bout en bout les percentiles p50, p99 et p99.9
à partir d'échantillons conjoints, par message.
\iqlookfor{horodatages externes et internes, joints par message, sous une charge réaliste.}
\end{solution}
